% Generated draft: requires mathematical and source review.
\documentclass[11pt]{article}
\usepackage[margin=25mm]{geometry}
\usepackage{fontspec}
\setmainfont{Latin Modern Roman}
\usepackage{amsmath,amssymb,mathtools,unicode-math}
\setmathfont{Latin Modern Math}
\usepackage{enumitem,needspace,xurl,hyperref}
\setlength{\parindent}{0pt}
\setlength{\parskip}{5pt}
\setlength{\emergencystretch}{3em}
\newcommand{\N}{\mathbb{N}}
\newcommand{\Z}{\mathbb{Z}}
\newcommand{\Q}{\mathbb{Q}}
\newcommand{\R}{\mathbb{R}}
\newcommand{\C}{\mathbb{C}}
\newcommand{\F}{\mathbb{F}}
\newcommand{\A}{\mathbb{A}}
\newcommand{\PP}{\mathbb{P}}
\newcommand{\E}{\mathbb{E}}
\newcommand{\eps}{\varepsilon}
\newcommand{\dd}{\,\mathrm{d}}
\begin{document}
\section*{UnsolvedMath ID: 20003150}
\textbf{Source problem code:} catalogue entry
\subsection{Definitions and mathematical statement}
Let $P\to B$ be a principal bundle for a compact connected Lie group $G$, with an $\operatorname{Ad}(G)$-invariant inner product on $\mathfrak g$. Work over a positive-dimensional base, $\dim B=2m\geq2$. A connection $\theta$ with curvature $\Omega$ is fat in the full principal-bundle sense if, for every nonzero $u\in\mathfrak g$, the horizontal two-form $\langle\Omega,u\rangle$ is nondegenerate at every point. Nondegenerate means symplectic as a bilinear form on the tangent space; closedness of each fiberwise two-form is not asserted in this terminology.

The canonical question asks whether the structure group of a fat principal $G$-bundle can be reduced. For a proper closed subgroup $H<G$, an underlying $H$-reduction is a principal $H$-subbundle $Q\subset P$, equivalently $P\cong Q\times_HG$. A reduction compatible with the given connection additionally requires that connection to preserve $Q$. These are different questions, and the source does not say which is intended or specify the base/group hypotheses.

For the connection-compatible reading, the aid reports a complete obstruction. Fatness makes the dual curvature map $\operatorname{ad}P\to\Lambda^2T^*B$ injective, hence the curvature map $\Lambda^2TB\to\operatorname{ad}P$ is pointwise surjective. Ambrose--Singer then gives full restricted holonomy $G$, so no proper $H$-reduction is preserved by that connection. This does not exclude a reduction of the underlying bundle with some other connection.

For any underlying $H$-reduction, write $\mathfrak h^\perp$ for the orthogonal complement of its Lie algebra and set
\[
 E_Q=Q\times_H\mathfrak h^\perp.
\]
The supplied aid reports that curvature gives a bundle monomorphism
\[
 E_Q\hookrightarrow\Lambda^2T^*B
\]
whose every nonzero fiber element is nondegenerate. This is a fiberwise linear symplectic pencil. For each $0\ne v\in T_bB$, contraction sends that pencil injectively into the annihilator of $v$, yielding the necessary bound
\[
 \dim(G/H)\leq\dim B-1.
\]
The reduction need not be preserved by the fat connection for this bundle-theoretic restriction. Conversely, a bundle monomorphism or the dimension inequality alone is not asserted sufficient to produce a reduction or a fat connection.

The aid further reports that proper underlying reductions are excluded for fat circle bundles over closed positive-dimensional bases and for fat $SU(2)$ bundles over $S^4$. The closed-base condition in the circle case is retained; it is not silently extended to open bases. The summary leaves the arbitrary nonnormal topological-reduction problem unresolved in its checked literature. No full classification of reductions or current literature assessment is supplied by this draft. The neighboring Problem 24 asks whether the group of a fat principal bundle must be $S^1$, $S^3$ or $SO(3)$; this contextual question is not treated as an assumed theorem.
\subsection{Short English statement}
A fat connection cannot preserve a proper reduction because its reported holonomy is full. Reduction of the underlying bundle is a separate question: the aid supplies a symplectic-pencil obstruction and two special exclusions, while leaving the general nonnormal case open.
\subsection{Sources}
\begin{itemize}
\item Local UnsolvedMath registry, permanent ID 20003150, source code catalogue entry; \texttt{aim-topology-notes.json}, notes[237].
\item AIM, \emph{Manifolds with nonnegative sectional curvature}, ``Group Actions and Submersions,'' Problem 25: \url{https://aimath.org/WWN/nnsectcurvature/nnsectcurvature.pdf}. The supplied verification reports the wording checked in the PDF and gives preceding Problem 24 as context; this draft does not independently repeat the check.
\item Ambrose--Singer holonomy and the reported curvature-pencil, circle-bundle and $SU(2)$-bundle obstructions are invoked in the aid without full bibliographic entries or exact theorem statements.
\item Machine-generated research attempt 1: partial result, partially solved, rigor labelled ``proved,'' low novelty confidence. Not peer reviewed; independently verify before citation. Registry updated 14 August 2026.
\end{itemize}
\end{document}
